% formal/sys-landscape/rotated-regular-products.tex — input by formal/main.tex
%
% Lagrangian products of rotated 2D polygons: systematic search for
% Viterbo counterexamples in the space of Lagrangian products.
%
% Sources:
%   - Generated data: experiments/sys-landscape/rotated-regular-products/lagrangian-products-5x5.jsonl
%     and experiments/sys-landscape/rotated-regular-products/lagrangian-products-<n>x<m>-6deg.jsonl
%   - HK-O 2024 counterexample: Haim-Kislev & Ostrover, Annals 2024
%   - Symmetry analysis: agent-written, Jörn-confirmed direction
%
% Structure:
%   - Mathematical setup (Lagrangian products, rotation parameter)
%   - Symmetry lemma (fundamental domain)
%
% Empirical content (sweep descriptions, figures, observations) lives in research/sys-landscape/design/rotated-regular-products.md.

\subsection{Lagrangian Products of Rotated Polygons}
\label{sec:lagrangian-products}

The only known 4D counterexample to Viterbo's conjecture is a Lagrangian
product of two pentagons \cite{HaimKislevOstrover2024}.
Our sample of random 4-polytopes never achieved~$\mathrm{sys} > 1$.
% See Section~\ref{sec:random-product-sample} for the random product sample experiment.

We now probe the conjecture by searching systematically in the space of Lagrangian products, where counterexamples are known to exist.

\subsubsection*{Mathematical setup}

Recall (Definition~\ref{def:lagrangian-product}) that the Lagrangian product
of convex polygons~$P \subset \mathbb{R}^2_q$ and~$Q \subset \mathbb{R}^2_p$
is $P \times_L Q = \{(q,p) : q \in P,\, p \in Q\} \subset \mathbb{R}^4$.

The 4-volume satisfies~$\operatorname{vol}_4(P \times_L Q)
= \operatorname{area}(P) \cdot \operatorname{area}(Q)$ (Fubini's theorem on
complementary subspaces).

\paragraph{Rotation parameter.}
Fix convex polygons~$P$ in $q$-space and~$Q$ in $p$-space.
For an angle~$\theta$, let~$R(\theta)$ denote
the standard rotation in~$\mathbb{R}^2$ and
define~$Q_\theta = R(\theta)\,Q$. We study the one-parameter family
\[
  K_\theta \;=\; P \times_L Q_\theta.
\]
Since rotating both factors by the same angle~$\varphi$ is a symplectic
transformation~$\operatorname{diag}(R(\varphi), R(\varphi)) \in \mathrm{Sp}(4)$,
only the \emph{relative} rotation between the factors affects~$\mathrm{sys}$.

\begin{unverified}
\begin{lemma}[Fundamental domain for rotation sweeps]
\label{lem:rotation-fundamental-domain}
Let~$P$ be a regular $n$-gon and~$Q$ a regular $m$-gon, both centered at the
origin. Then~$\mathrm{sys}(P \times_L Q_\theta)$ depends on~$\theta$ only
through its equivalence class in~$[0, \pi/\operatorname{lcm}(n,m)]$.
\end{lemma}

% [TODO: JÖRN - verify lemma proof below]
\begin{proof}
\emph{Periodicity.}
The block-diagonal matrix~$M_\varphi = \operatorname{diag}(R(\varphi),
R(\varphi))$ is symplectic (direct verification: $M_\varphi^T J_0 M_\varphi
= J_0$). Applying~$M_{2\pi/n}$ to~$P \times_L Q_\theta$
maps~$P$ to itself ($n$-fold rotational symmetry)
and replaces~$Q_\theta$ with~$Q_{\theta + 2\pi/n}$. Similarly, the $m$-fold symmetry of~$Q$
gives $Q_\theta = Q_{\theta + 2\pi/m}$. A function with periods~$a$
and~$b$ also has period~$\gcd(a,b)$; since
$\gcd(2\pi/n,\, 2\pi/m) = 2\pi/\operatorname{lcm}(n,m)$,
the combined periodicity is~$2\pi/\operatorname{lcm}(n,m)$.

\emph{Mirror symmetry.}
The matrix~$M = \operatorname{diag}(-1, 1, -1, 1)$ is symplectic
(verification:~$M^T J_0 M = J_0$). It acts on~$\mathbb{R}^2_q$ as reflection
across the~$q_2$-axis and on~$\mathbb{R}^2_p$ as reflection across
the~$p_2$-axis. For our angle convention (normals starting at~$\pi/2$), both
regular polygons are invariant under these reflections, and~$M$
maps~$Q_\theta$ to~$Q_{-\theta}$.

Together: $\mathrm{sys}(\theta)$ has period~$2\pi/\operatorname{lcm}(n,m)$ and
satisfies~$\mathrm{sys}(\theta) = \mathrm{sys}(-\theta)$, so the fundamental
domain is~$[0, \pi/\operatorname{lcm}(n,m)]$.
\end{proof}
\end{unverified}

% Factor swap: N = -J_0 is symplectic, maps P x_L Q to Q x_L (-P).
% For centrally symmetric polygons, swapping factors preserves sys.
% For odd polygons, -P = R(π)P so the swap also rotates one factor by 180°.
% (Jörn verified 2025-02-14.)
